\section{Pre-registration, deviations and reproduction}
\label{app:reproduce}

\subsection{What was fixed, and when}

The protocol was not written in one sitting, and describing it as a single
pre-registration would overstate it. The sequence was:

Note first that the configuration numbers in this paper are not chronological.
Config~2 ran first; Configs~1 and~3 are the second task.

\begin{enumerate}
\item A single-seed \emph{smoke test} was frozen on 2026-07-19 and run on
CodeFinQA. Its decision quantity was GRPO minus matched SFT---one opponent, not
a $\max$---with a Go bar at five points and a No-Go floor at three. Its own text
states that single-seed results are smoke evidence and not paper evidence, and
no number from it appears here.
\item The protocol this paper reports was written on 2026-07-20 as a new
question rather than a patch on the first. It introduced the five arms, the
$\max$ aggregation over $\mathcal{B}$, the $+3$ Go criterion with its companion
conditions, the operating-band rule for $m_L$ and $m_H$, and the good-faith
clause forbidding us from weakening the solver if GRPO won. It was written for
CodeFinQA, which this paper calls \textbf{Config~2}.
\item Config~2 was executed.
\item An addendum was frozen on 2026-07-24, \emph{before any CodeTAT-QA
training or evaluation run}, instantiating CodeTAT-QA as the second task and
carrying the Gate~0/1/2 criteria, the definition of $\Delta$ and the verdict
vocabulary over verbatim.
\item Configs~1 and~3 (CodeTAT-QA at 0.5B and 7B) were executed under it.
\end{enumerate}

The honest labelling that follows is the reverse of what the numbering
suggests. \textbf{Config~2 is the exploratory one}: its protocol was written by
authors who had already seen single-seed results on that same dataset.
\textbf{Configs~1 and~3 are confirmatory}: the rule governing them was frozen
before any of their data existed. The abstract's framing---``pre-registered,
before the second task was run''---refers to step 4 and is accurate, but a
reader could take it to mean the whole study was registered before any data
existed, and it was not. Note also that the quantitative core of this paper, the
11 groups at 0.5B and 1B, straddles the two: six of them are confirmatory and
five are not.

One further qualification on the word ``frozen.'' The document written in step 2
carries a status line reading \emph{not yet frozen}, and the addendum of step 4
refers to the operative protocol as that document \emph{together with} a
numbered entry in the deviation log. The frozen protocol is therefore a draft
plus one logged amendment rather than a single signed artifact, and a reader
reconstructing it must read both. We publish both.

Within the confirmatory scope, only ``GRPO vs.\ the strongest non-RL arm'' is a
confirmatory comparison; the ablations, audits, sweep, dose curves and every
analysis in Section~\ref{sec:statistics} of this extended version are
exploratory and are reported as such.

One rule was adopted after the fact: the exclusion of a run whose positive
reward rate is identically zero. It was not specified in advance. Its effect is
reported in both directions in Appendix~\ref{app:stability}.

\subsection{Artifacts}

The following are released with this extended version, at
\url{\repourl}:

\begin{itemize}
\item \textbf{The pre-registration documents and their addendum}, verbatim, in
the state they were frozen.
\item \textbf{The deviation logs}, listing every procedural departure from the
protocol and the direction in which it influences $\Delta$. These are not
curated: they include the incident in which an early verdict script silently
omitted the solver from the $\max$ set and reported $\Delta = +3.00$ where the
corrected value is $-14.00$. That error is why the aggregation is computed over
an explicit arm list and checked, and we regard publishing it as part of the
claim.
\item \textbf{The solver development logs}, round by round, with the dev
accuracy after each round and the failure mode each round targeted
(Section~\ref{subsec:solvercost}).
\item \textbf{The solver, the verifier and the training/evaluation harness.}
\item \textbf{Per-item outcome vectors} for all five arms in all 17 groups: the
inputs to every number in Section~\ref{sec:statistics} and
Appendix~\ref{app:allarms}.
\item \textbf{The inference-engine equivalence gate} of
Section~\ref{subsec:enginegate}: its script, its raw result and its 18 per-item
flip samples.
\item \textbf{The analysis scripts.} Every table in this document that reports
a statistic is emitted directly from the stored artifacts by
\texttt{gen\_extended\_\allowbreak stats.py} or
\texttt{gen\_supp\_\allowbreak tables.py} under \texttt{analysis/}; both carry
self-checks that abort if a recomputed value disagrees with the value printed
in the paper.
\end{itemize}

\subsection{The inference-engine equivalence gate}
\label{subsec:enginegate}

One entry in the deviation log bears directly on whether a $3$-point Go
threshold is measurable at all, so we state it here rather than leaving it to a
reader of the logs.

Early in the study we considered replacing HuggingFace \texttt{generate} with
vLLM, which is roughly $60\times$ faster on this workload. Before adopting it we
ran an equivalence gate: the same checkpoint, the same 200 dev items, a prompt
fingerprint verified identical on both sides, and---by design---generation
separated from verification, so that the two engines only produced text and a
single shared verifier scored it. Anything else would have folded environment
differences into the measurement.

\begin{center}
\begin{tabular}{lr}
\toprule
HuggingFace \texttt{generate} accuracy & $41.00\%$ \\
vLLM accuracy & $38.00\%$ \\
difference & $-3.00$ pp \\
\midrule
per-item reward agreement & $91.00\%$ (18/200 flipped) \\
generated text identical & $38.50\%$ \\
HF-only correct / vLLM-only correct & 12 / 6 \\
McNemar $\chi^2$ & $1.39$ (n.s.) \\
\bottomrule
\end{tabular}
\end{center}

The engine alone moves accuracy by $3.00$ points, which is exactly this study's
pre-registered Go threshold, and it does so in no systematic direction. A
constant offset could have been corrected; non-systematic noise of the same
magnitude as the effect under test cannot. The mechanism is not floating-point
jitter at the last digit but path divergence: a 0.5B policy reading a
${\sim}3000$-token prompt in bf16 has very small logit gaps, so a single
\texttt{argmax} flip early in decoding sends the emitted program down a
different branch. Only $38.5\%$ of completions are textually identical, which
shows the divergence is the normal case rather than an occasional accident.

We therefore rejected the substitution and used HuggingFace \texttt{generate}
throughout. We report this for two reasons. First, it is the reason a reader
should believe the $3$-point threshold means anything in this study: we checked
that our own instrument does not move by that much for reasons unrelated to the
treatment. Second, it generalizes the paper's argument one level down. The
headline result is that a reported RLVR effect depends on which control arm you
compare against, to the tune of $19$ points; this gate shows that it also depends
on which inference engine you decode with, to the tune of $3$---and that an
author who swapped engines mid-study would see neither dependence in the
reported numbers.

\subsection{Reproducing the numbers}

With the artifact bundle in place:

\begin{verbatim}
# results_backup/ -> the fact table
python analysis/gen_data_facts.py
# ablation, PoT, audit, stability
python analysis/gen_supp_tables.py
# per-group stats, aggregators, cross-fitting,
# coverage, zero-advantage
python analysis/gen_extended_stats.py
\end{verbatim}

Randomized quantities are seeded: the bootstrap uses $B = 10{,}000$ resamples
and the cross-fitted estimator $K = 2{,}000$ splits, both under a fixed seed, so
the reported intervals are reproducible to the digit.

\subsection{Hardware and software}

One NVIDIA RTX A6000 (48\,GB) per run; PyTorch 2.10, \texttt{transformers}
4.57, \texttt{peft} 0.19, \texttt{trl} 1.8. Every arm trains a LoRA adapter
($r{=}16$, $\alpha{=}32$) at learning rate $10^{-5}$; GRPO uses 8 generations
per prompt, $\beta{=}0.04$, temperature $1.0$, top-$p$ $0.95$ and effective
batch 64, reduced to per-device 4 with accumulation 16 at 7B. The three arms
that keep training are each capped at 5700\,s of wall clock. The solver runs on
CPU in seconds.
